\documentclass[a4paper,11pt]{article}
\usepackage[utf8]{inputenc}
\usepackage[T1]{fontenc}
\usepackage[french]{babel}
\usepackage[left=2cm,right=2cm,top=2cm,bottom=2cm]{geometry}
\usepackage{xcolor}
\usepackage{hyperref}
\usepackage{fontawesome5}
\usepackage{parskip}

% --- Couleurs (Bleu LCL / Crédit Agricole) ---
\definecolor{primary}{RGB}{0, 51, 102}
\hypersetup{colorlinks=true, urlcolor=primary}

\begin{document}

% --- EXPÉDITEUR ---
\noindent
\begin{minipage}[t]{0.5\textwidth}
    \textbf{Loïc Aron MBASSI EWOLO}\\
    Élève-Ingénieur Bac+5 -- Double diplôme ENSEA\\[3pt]
    \faMapMarker\ Cergy, France\\
    \faPhone\ (+33) 7 59 52 33 98\\
    \faEnvelope\ \href{mailto:loicaron.mbassiewolo@ensea.fr}{loicaron.mbassiewolo@ensea.fr}
\end{minipage}
\hfill
% --- DATE ---
\begin{minipage}[t]{0.4\textwidth}
    \raggedleft
    Cergy, le 27 février 2026
\end{minipage}

\vspace{1.2cm}

% --- DESTINATAIRE ---
\hfill
\begin{minipage}[t]{0.5\textwidth}
    \raggedleft
    \textbf{LCL -- Direction des Systèmes d'Information}\\
    Service Recrutement\\
    Villejuif, France
\end{minipage}

\vspace{1cm}

\noindent\textbf{Objet :} Candidature en alternance -- Ingénieur d'Études et Développement Full Stack

\vspace{0.8cm}

Madame, Monsieur,

Trois ans de développement Full Stack en contexte professionnel, dont une expérience en secteur Fintech/bancaire, m'ont permis de construire une maîtrise concrète du cycle de vie des applications d'entreprise : conception, développement, tests et maintenance évolutive. En double diplôme à l'ENSEA (Systèmes Embarqués, IA) après un diplôme d'ingénieur à l'École Polytechnique de Yaoundé, je souhaite rejoindre la tribu CRM \& Communication client de LCL en alternance pour contribuer au développement des outils qui accompagnent vos conseillers au quotidien.

Chez CSC (Cameroon Software Company), dans le secteur Fintech, j'ai conçu le backend d'une plateforme de transactions financières sécurisée en Laravel, avec gestion de transactions concurrentes et conformité aux normes du secteur. Ce contexte bancaire m'a familiarisé avec les exigences de fiabilité, de sécurité et de traçabilité propres aux SI financiers. En parallèle, mon expérience de 2 ans chez SOSDOCTA (télémédecine) m'a confronté à la maintenance évolutive d'une application critique, avec gestion de données sensibles (RGPD), sécurisation OAuth2/JWT et travail en remote avec une équipe internationale.

Côté Java, j'ai développé une application covoiturage avec Spring Boot et ScyllaDB lors du projet Freengo, ainsi qu'une application Desktop en Java utilisant intensivement les Design Patterns (Singleton, Factory, Observer, MVC). Mon expérience frontend en React.js et TypeScript (marketplace TOGETTECH, projet ASSO en stack MERN) me permettra une montée en compétence rapide sur Angular, les fondamentaux du développement composant, du binding de données et de la gestion d'état étant transversaux. J'accorde une attention particulière aux tests unitaires et d'intégration, à la qualité du code et aux revues en équipe.

Je suis disponible dès septembre 2026 pour une alternance d'un an. Intégrer la DSI de LCL représente pour moi l'opportunité de travailler sur des applications Java/Angular à fort impact métier, dans un environnement technologique exigeant (Cloud, Big Data, IA) et au sein d'une équipe structurée avec l'accompagnement de vos tech leads.

Je reste à votre disposition pour un entretien afin de discuter de ma candidature.

Je vous prie d'agréer, Madame, Monsieur, l'expression de mes salutations distinguées.

\vspace{0.8cm}

\textbf{Loïc Aron MBASSI EWOLO}\\
\textit{Élève-Ingénieur ENSEA / Polytechnique Yaoundé}

\end{document}
